\SourceLeaf{226}%
% A544-P0226-D01
\[
PA+QB\equiv1
\]
Porro vero manifestum est,
\[
\begin{aligned}
\text{Dimens. ipsius }B+\text{Dim. ipsius }a&\text{ esse }=\text{Dim. }A.\\
\text{Dim. }C+\text{Dim. }b&=\pdot\text{Dim. }B\\
&\text{etc.}\\
\text{Dim. }L+\text{Dim. }k&=\text{Dim. }K
\end{aligned}
\]
Quare
\[
\text{Dim. }L+\text{Sum. Dim. }a,b,\dotstwo k=\text{Dim. }A
\]
Patet vero dimensionem ipsius $a^{(\kappa)}$ adeoque etiam
\[
\text{Dim. ipsius }Q\text{ esse }=\text{Sum. Dim. }a,b,c,\dotstwo\text{ i. e. }=\alpha-\text{Dim. }L
\]
itaque
\[
\text{Dim. ipsius }P=\beta-\text{Dim. }L\qquad\text{Q. E. D.}
\]

% A544-P0226-D02
\Article{336}
Hinc autem sequitur, si $M$ est divisor communis maximae dimensionis functionum $A,B$, semper poni posse
\[
AP+BQ\equiv M
\]

% A544-P0226-D03
Exempla praecedentis theorematis brevitatis gratia omitto, sed lectores non negligent, per ea facilitatem huius generis problemata tractandi sibi comparare. Ceterum operae pretium erit admonere, theorema praecedens etiam de functionibus absolute sumtis valere, quarum quidem coëfficientes sint numeri rationales. Hoc ex demonstrationis modo per se elucebit. Nobis autem ei rei immorari non licet. Similia lector etiam non admonitus in sequentibus observabit.

% A544-P0226-D04
Si $A$ nec cum $B$ nec cum $C$ divisorem ullius dimensionis communem habeat, etiam cum producto $BC$ nullum habebit divisorem communem. Sit enim
\[
PA+QB\equiv1,\quad\text{erit}\quad PAC+QBC\equiv C
\]
Iam si $A$ cum $BC$ divisorem $M$ communem haberet, hic etiam ipsam $C$ divideret contra hyp. Hinc generaliter si functio $A$ ad $B,C,D$ etc. prima, etiam ad omnium productum erit prima.

% A544-P0226-D05
Si $A,B,C,D$ etc. nullum divisorem habeant omnibus communem, fieri potest
\SourceLeaf{227}%
% A544-P0227-D01
\[
PA+QB+RC+SD+\text{etc.}\equiv1
\]
Sit divisor maximae dimensionis inter $A$ et $B,M$; inter $M$ et $C,M'$; inter $M'$ et $D,M''$ etc.: patet, ultimum huius seriei terminum fore nullius dimensionis (hyp.). Quare poni poterit
\[
aA+bB\equiv M,\quad mM+cC\equiv M',\quad m'M'+dD\equiv M'',\quad\text{etc.}
\]
unde substitutionibus factis theorematis veritas apparet.

% A544-P0227-D02
\Article{337}
\textsc{Theorema.} \emph{Si $A,B,C$ etc. sint functiones inter se primae (quarum binae quaeque nullum habeant divisorem communem) secundum modulum $p$, et functio $M$ secundum eundem modulum per singulas sit divisibilis; etiam per omnium productum erit divisibilis.}

% A544-P0227-D03
\emph{Demonstr.} Poni enim potest $PA+QB\equiv1$, quare erit
\[
\frac MA Q+\frac MB P\equiv\frac M{AB}
\]
Iam quum $C$ ad $AB$ prima, erit etiam $M$ per $ABC$ divisibilis similique ratiocinio per $ABCD$ etc.

% A544-P0227-D04
\Article{338}
Si congruentia $\xi\equiv0$ habeat radices $x\equiv a$, $x\equiv b$, $x\equiv c$ etc., $\xi$ per productum ex $(x-a)$, $(x-b)$, $(x-c)$ etc. dividi poterit; cum enim $a,b,c\pdot$ etc. supponantur incongrui, functiones $x-a,x-b,x-c$ etc. erunt primae inter se, et quum $\xi$ per singulas dividatur, etiam per productum ex omnibus dividetur. Hinc patet, radicum multitudinem congruentiae dimensionem superare non posse: quae est demonstratio huius theorematis, quam polliciti sumus.

% A544-P0227-D05
Sed simul hinc perspicitur, quomodo congruentiarum solutio partem tantummodo constituat multo altioris disquisitionis, scilicet de resolutione functionum in factores. Manifestum est, congruentiam $\xi\equiv0$ nullas habere radices reales, si $\xi$ nullos factores unius dimensionis habeat; at hinc nihil obstat, quominus $\xi$ in factores duarum, trium pluriumve dimensionum resolvi possit, unde radices quasi \emph{imaginariae} illi attribui possint. Revera, si simili licentia, quam recentiores mathematici usurparunt, uti talesque quantitates imaginarias introducere 
\SourceLeaf{228}%
% A544-P0228-D01
voluissemus, omnes nostras disquisitiones sequentes incomparabiliter contrahere licuisset; sed nihilominus maluimus omnia ex principiis deducere\begingroup\renewcommand{\thefootnote}{*)}\footnote{Alia forsan occasione de hac re opinionem nostram fusius explicabimus.}\endgroup.

% A544-P0228-D02
\Article{339}
\emph{Functiones} secundum modulum determinatum \emph{primae} vocantur, quae per nullas functiones inferiorum dimensionum secundum hunc modulum dividi possunt.

% A544-P0228-D03
Ita omnes functiones unius dimensionis erunt primae, functiones autem duarum dimensionum aut erunt primae aut ex binis unius dimensionis compositae: quare $\xi$ erit functio prima duarum dimensionum, si congruentia $\xi\equiv0$ nullas radices reales admittit. Ex. gr. $xx+x+1$ pro modulo $5$ est prima, quia
\[
xx+x+1\equiv(x-2)^2-3\;(\mathrm{mod.}\,5)
\]
et $3$ non-residuum quadraticum numeri $5$.

% A544-P0228-D04
Hae vero functiones primae prae omnibus attentionem nostram desiderant. Quamvis enim aliae quam primi gradus ad inveniendas radices reales inservire non possint, amplior earum consideratio tum ob insignes ipsarum proprietates tum ob alias egregias veritates ex his deducendas sese commendat.

% A544-P0228-D05
\Article{340}
\textsc{Theorema.} \emph{Functio quaecunque aut est prima aut ex functionibus primis composita; posteriorique casu unico tantum modo e functionibus primis componi potest.}

% A544-P0228-D06
\emph{Demonstr.} Nisi enim functio proposita $A$ sit prima, per aliam inferioris dimensionis $B$ dividetur. Si $B$ non est functio prima, per aliam $C$ inferioris gradus dividetur, itaque pergendo patet, tandem ad functionem primam deveniri, quoniam alias haec series foret infinita, quod, quoniam dimensiones perpetuo decrescunt, absurdum est. Iam si ultima functio prima sit $L$, haec omnes antecedentes metietur. Quare $A\equiv LA'$ eritque $A'$ inferioris dimensionis quam $A$. Quod iterum fiet $A'\equiv L'A''$ etc., patet, tandem ad functionem primam perveniri, adeoque $A$ erit $\equiv$ producto e functionibus primis $L,L',L''$ etc.\ $\pdot$Q.\,E.\,Pr.

% A544-P0228-D07
Iam si etiam esset $A\equiv MM'M''$ etc. neque omnes $L,L',L''$ etc. eaedem cum omnibus $M,M',M''$ etc., eiiciantur eae, quae utrique seriei communes
\SourceLeaf{229}%
% A544-P0229-D01
sunt. Remaneantque $\lambda,\lambda',\lambda''\dotstwo$; $\mu,\mu',\mu'',\ldots$ eritque $\mu$ ad $\lambda,\lambda',\lambda''$ etc. prima, quare etiam ad productum $\lambda\lambda'\lambda''$ etc.; tamen esse debet
\[
\lambda\lambda'\lambda''\dotstwo\equiv\mu\mu'\mu''\dotstwo\text{ i. e. }\frac{\lambda\lambda'\lambda''\dotstwo}{\mu}\equiv\mu'\mu''\dotstwo\quad\text{Q. E. A.}
\]

% A544-P0229-D02
\Article{341}
Primum caput harum investigationum in eo consistet, ut functionum primarum cuiusvis dimensionis multitudinem determinemus. Quoniam enim pro modulo determinato numerus omnium functionum diversarum (incongruarum) cuiuslibet gradus est definitus, ex his vero aliae sunt ex primis inferiorum graduum compositae, aliae primae, etiam harum numerus finitus erit. Rigorosa huius rei evolutio satis est lubrica; a casibus simplicioribus incipiemus.

% A544-P0229-D03
Posito modulo $=p$, numerus omnium functionum diversarum $n^{\mathrm{ti}}$ gradus huius formae
\[
x^n+Ax^{n-1}+Bx^{n-2}+Cx^{n-3}+\text{etc.}
\]
erit $p^n$; coëfficientium enim $\pdot A,B\pdot C$ etc. numerus est $n$; et quum quivis independenter a reliquis possit esse $\equiv0,1,2,3\dotstwo(p-1)(\mathrm{mod.}\,p)$, ex combinationum theoria sequitur, $p^n$ combinationes diversas haberi, quae igitur omnium functionum diversarum huius gradus complexum definiunt.

% A544-P0229-D04
Ita functiones unius dimensionis erunt $p$, scilicet $x,x+1,x+2$ usque ad $x+p-1$; functiones duarum dimensionum $pp$ etc.

% A544-P0229-D05
\Article{342}
Iam supra monuimus, omnes functiones primi gradus pro primis habendas esse; si igitur, quod ad propositum nostrum sufficit, ad eas functiones nos restringamus, quarum terminus summus habet coëfficientem $1$, erunt $p$ functiones primi gradus seu unius dimensionis.

% A544-P0229-D06
Functiones secundi gradus omnes aut e binis primi gradus erunt compositae aut primae. Iam ex combinationum theoria constat, $p$ res diversas admissis repetitionibus $\frac{p\pdot p+1}{1\pdot2}$ modis diversis combinari posse, quare totidem functiones erunt e binis primis unius dimensionis compositae, adeoque $pp-\frac{p\pdot p+1}{1\pdot2}=\frac12(pp-p)$ functiones primae duarum dimensionum.

\SourceLeaf{230}%
% A544-P0230-D01
Simili modo e functionibus omnibus tertii gradus, quarum numerus est $p^3$, excludendae sunt eae, quae e ternis primis unius dimensionis componuntur, quarum numerus est $\frac{p\pdot p+1\pdot p+2}{1\pdot2\pdot3}$; insuperque eae, quae e functione prima unius aliaque duarum dimensionum componuntur, quarum numerus est $p\pdot\frac12(pp-p)$; quibus deletis restabunt $\frac13(p^3-p)$; tot igitur sunt primae trium dimensionum. Elucet hoc modo semper continuari posse.

% A544-P0230-D02
\Article{343}
Ut autem hae operationes facilius absolvantur simulque ad evolutionem legis generalis via sternatur, rem generaliter considerabimus. Brevitatis gratia designamus per $(1)$ multitudinem functionum primarum unius dimensionis, per $(2)$ numerum functionum primarum duarum dimensionum, sic porro per $(1^2)$ multitudinem functionum e binis primis unius dimensionis compositarum etc. etc., generaliter per $(1^\alpha2^\beta3^\gamma\ldots)$ multitudinem functionum omnium, quae e functionibus primis compositae sunt, scilicet ex $\alpha$ unius, $\beta$ duarum, $\gamma$ trium etc. dimensionum, quarum itaque dimensio erit $\alpha+2\beta+3\gamma+\text{etc.}$ Tum per praecedentia theoriamque combinationum elucet, fore
\[
(1^\alpha2^\beta3^\gamma4^\delta\ldots)=(1^\alpha)(2^\beta)(3^\gamma)(4^\delta)\dotstwo
\]
\[
(1^\alpha)=\frac{(1)\pdot(1)+1\pdot(1)+2\pdot(1)+3\ldots(1)+\alpha-1}{1\pdot2\pdot3\pdot4\ldots\alpha}
\]
seu generaliter
\[
(a^\alpha)=\frac{(a)\pdot(a)+1\pdot(a)+2\pdot(a)+3\ldots(a)+\alpha-1}{1\pdot2\pdot3\pdot4\ldots\alpha}
\]

% A544-P0230-D03
Denique manifestum est, si omnes modi diversi numerum $n$ e numeris $1,2,3,\ldots$ per additionem componendi colligantur, qui designentur per $\alpha\pdot1+\beta\pdot2+\gamma\pdot3+\text{etc.}$, summam omnium harum expressionum $(1^\alpha2^\beta3^\gamma\dotstwo)$ aequalem fore multitudini omnium functionum $n$ dimensionum, i. e. $=p^n$. Ita
\[
\begin{aligned}
p&=(1)\\
pp&=(1^2)+(2)\\
p^3&=(1^3)+(1\pdot2)+(3)\\
p^4&=(1^4)+(1^2\pdot2)+(1\pdot3)+(2^2)+(4)\\
&\text{etc.}
\end{aligned}
\]

% A544-P0230-D04
Perspicuum est, in expressione $p^n$ praeter quantitates $(1),(2),(3)$ etc. etiam hanc
